
% ============================================================
% FRENCH TENSES — FUTUR ANTÉRIEUR
% Version graphique LaTeX/TikZ
% Compilation : XeLaTeX
% ============================================================
\documentclass[a4paper,10pt]{article}

\usepackage[a4paper,top=9mm,bottom=8mm,left=11mm,right=11mm]{geometry}
\usepackage{fontspec}
\usepackage[french]{babel}
\usepackage{microtype}
\usepackage{tikz}
\usetikzlibrary{arrows.meta,calc,positioning}
\usepackage[most]{tcolorbox}
\usepackage{tabularx}
\usepackage{array}
\usepackage{xcolor}
\usepackage[normalem]{ulem}

\setmainfont{Noto Sans}
\setsansfont{Noto Sans}

% ------------------------------------------------------------
% PALETTE — FUTUR = ORANGE
% ------------------------------------------------------------
\definecolor{Orange}{HTML}{D8842F}
\definecolor{OrangeDark}{HTML}{9A5716}
\definecolor{OrangeLight}{HTML}{FFF1DF}
\definecolor{OrangeLine}{HTML}{E7B77E}
\definecolor{Border}{HTML}{E7C7A6}

\definecolor{Red}{HTML}{D52F35}
\definecolor{RedLight}{HTML}{FCEBED}
\definecolor{Gold}{HTML}{B27A00}
\definecolor{GoldLight}{HTML}{FFF3D7}
\definecolor{Ink}{HTML}{172A32}
\definecolor{Grey}{HTML}{64757D}
\definecolor{White}{HTML}{FFFFFF}

\pagestyle{empty}
\setlength{\parindent}{0pt}
\setlength{\parskip}{0pt}
\setlength{\tabcolsep}{3pt}

% ------------------------------------------------------------
% CADRE
% ------------------------------------------------------------
\newcommand{\PageFrame}{%
\begin{tikzpicture}[remember picture,overlay]
  \coordinate (A) at ([xshift=8mm,yshift=-8mm]current page.north west);
  \coordinate (B) at ([xshift=-8mm,yshift=-8mm]current page.north east);
  \coordinate (C) at ([xshift=-8mm,yshift=8mm]current page.south east);
  \coordinate (D) at ([xshift=8mm,yshift=8mm]current page.south west);
  \def\r{1.7mm}

  \draw[Border,line width=.65pt]
    (A)
    -- ([xshift=-\r]B)
    arc[start angle=90,end angle=0,radius=\r]
    -- ([yshift=\r]C)
    arc[start angle=0,end angle=-90,radius=\r]
    -- (D) -- cycle;

  \draw[Orange,line width=3.2mm] (A) -- (D);
\end{tikzpicture}%
}

% ------------------------------------------------------------
% BOÎTES
% ------------------------------------------------------------
\tcbset{
  enhanced,
  boxrule=.6pt,
  arc=2.8mm,
  left=3.7mm,right=3.7mm,top=1.75mm,bottom=1.75mm,
  boxsep=1mm,
  coltext=Ink
}

\newtcolorbox{bluebox}{colback=OrangeLight,colframe=OrangeLine}
\newtcolorbox{redbox}{colback=RedLight,colframe=Border}
\newtcolorbox{whitebox}{colback=White,colframe=Border}
\newtcolorbox{goldbox}{
  colback=GoldLight,colframe=GoldLight,
  boxrule=0pt,arc=2.5mm,
  left=2.7mm,right=2.7mm,top=1.25mm,bottom=1.25mm
}

\newcommand{\sectiontitle}[2]{%
  {\large\bfseries\textcolor{#1}{#2}}%
}

\newcommand{\pill}[2][OrangeLight]{%
  \tikz[baseline=(X.base)]{
    \node[
      fill=#1,rounded corners=2pt,
      inner xsep=4.2pt,inner ysep=1.6pt,
      font=\small\bfseries,text=Ink
    ] (X) {#2};
  }%
}

\newcommand{\signalpill}[1]{%
  \tikz[baseline=(X.base)]{
    \node[
      fill=GoldLight,rounded corners=2pt,
      inner xsep=3.7pt,inner ysep=1.1pt,
      font=\scriptsize\bfseries,text=Ink
    ] (X) {#1};
  }%
}

% ------------------------------------------------------------
% TIMELINE
% ------------------------------------------------------------
\newcommand{\Timeline}{%
\begin{tikzpicture}[x=1cm,y=1cm]
  \draw[-{Latex[length=3mm]},very thick,Ink] (0,0) -- (15.3,0);

  \draw[-{Latex[length=3mm,width=2mm]},very thick,Orange]
    (8.2,1.10) -- (8.2,.10);
  \node[font=\bfseries\small,text=Orange] at (8.2,1.32) {MAINTENANT};

  \draw[Orange,line width=2.2pt] (9.3,0) -- (11.8,0);
  \fill[Orange] (9.3,0) circle (3.2pt);
  \fill[Orange] (11.8,0) circle (3.2pt);

  \node[
    fill=OrangeLight,rounded corners=2pt,
    inner xsep=5pt,inner ysep=2pt,
    font=\scriptsize\bfseries,text=OrangeDark
  ] at (10.55,.72) {action terminée};

  \node[font=\scriptsize\bfseries,text=Grey] at (10.55,-.55)
    {avant un moment futur};

  \node[font=\scriptsize\bfseries,text=Ink] at (11.8,.34) {repère futur};

  \node[font=\scriptsize\bfseries,text=Grey] at (3.0,-.48) {PASSÉ};
  \node[font=\scriptsize\bfseries,text=Orange] at (8.2,-.48) {PRÉSENT};
  \node[font=\scriptsize\bfseries,text=Grey] at (13.7,-.48) {FUTUR};
\end{tikzpicture}%
}

\begin{document}
\PageFrame

% ------------------------------------------------------------
% TITRE
% ------------------------------------------------------------
\vspace*{1mm}
\begin{center}
  {\fontsize{25}{27}\selectfont\bfseries\textcolor{Ink}{FUTUR ANTÉRIEUR}}

  \vspace{0.75mm}
  {\large\textcolor{Grey}{Le futur antérieur}}

  \vspace{3mm}

  \begin{tcolorbox}[
    width=.88\textwidth,
    colback=OrangeLight,
    colframe=OrangeLight,
    boxrule=0pt,
    arc=3mm,
    halign=center,
    top=2.2mm,bottom=2.2mm
  ]
    {\large\bfseries\textcolor{OrangeDark}{AUXILIAIRE AU FUTUR SIMPLE + PARTICIPE PASSÉ}}
  \end{tcolorbox}
\end{center}

\vspace{1mm}

% ------------------------------------------------------------
% TIMELINE
% ------------------------------------------------------------
\begin{bluebox}
  \sectiontitle{Orange}{La ligne du temps}

  \vspace{0.75mm}
  \begin{center}
    \Timeline
  \end{center}

  \vspace{-1mm}
  \begin{center}
    \small
    Le \textbf{futur antérieur} présente une action
    \textbf{\textcolor{Orange}{qui sera terminée avant un autre moment futur}}.
  \end{center}
\end{bluebox}

\vspace{1.35mm}

% ------------------------------------------------------------
% CONSTRUCTION
% ------------------------------------------------------------
\begin{bluebox}
  \sectiontitle{Orange}{Construction}

  \vspace{0.75mm}
  \begin{tabularx}{\textwidth}{@{}p{.29\textwidth} X@{}}
    \textbf{Auxiliaire} &
      \textbf{avoir} ou \textbf{être} au futur simple.\\
    \textbf{Participe passé} &
      \textbf{j'aurai fini}, \textbf{nous serons partis}, \textbf{elle aura compris}.\\
    \textbf{Formule} &
      \textbf{AUXILIAIRE AU FUTUR SIMPLE + PARTICIPE PASSÉ}.
  \end{tabularx}

  \vspace{0.7mm}
  \textit{Avec \textbf{être}, le participe passé s'accorde avec le sujet.}
\end{bluebox}

\vspace{1.35mm}

% ------------------------------------------------------------
% EXEMPLES / ATTENTION
% ------------------------------------------------------------
\begin{minipage}[t]{.485\textwidth}
\begin{bluebox}
  \sectiontitle{Orange}{Exemples}

  \vspace{0.75mm}
  Quand tu arriveras, j'\textbf{aurai déjà fini}.\\[-1mm]
  Nous \textbf{serons partis} avant la nuit.\\[-1mm]
  Elle \textbf{aura terminé} son travail à 18 heures.
\end{bluebox}
\end{minipage}
\hfill
\begin{minipage}[t]{.485\textwidth}
\begin{redbox}
  \sectiontitle{Red}{Attention !}

  \vspace{0.75mm}
  Le futur antérieur exprime une action
  \textbf{achevée avant un repère futur}.

  \vspace{0.4mm}
  Ne pas le confondre avec le futur simple :
  \textbf{j'aurai fini} / \textbf{je finirai}.
\end{redbox}
\end{minipage}

\vspace{1.35mm}

% ------------------------------------------------------------
% USES
% ------------------------------------------------------------
\begin{bluebox}
  \sectiontitle{Orange}{Quand utilise-t-on le futur antérieur ?}

  \vspace{0.75mm}
  \begin{tabularx}{\textwidth}{@{}p{.31\textwidth} X@{}}
    \pill{ANTÉRIORITÉ} &
      Action terminée avant une autre action future.
      \newline \textit{Quand tu viendras, j'aurai terminé.}\\[0.7mm]

    \pill{RÉSULTAT FUTUR} &
      Résultat qui sera obtenu à un moment futur.
      \newline \textit{À midi, nous aurons fini.}\\[0.7mm]

    \pill{ÉCHÉANCE} &
      Action accomplie avant une limite ou une date future.
      \newline \textit{D'ici vendredi, elle aura répondu.}\\[0.7mm]

    \pill{CHRONOLOGIE} &
      Première action dans une succession d'événements futurs.
      \newline \textit{Quand il sera parti, nous commencerons.}
  \end{tabularx}
\end{bluebox}

\vspace{1.35mm}

% ------------------------------------------------------------
% COMPARISON
% ------------------------------------------------------------
\begin{whitebox}
  \sectiontitle{Orange}{Futur antérieur ou futur simple ?}

  \vspace{0.65mm}
  \begin{tabularx}{\textwidth}{@{}p{.45\textwidth} X@{}}
    \textbf{\textcolor{Orange}{Futur antérieur}} &
      \textbf{\textcolor{Grey}{Futur simple}}\\
    action terminée avant un repère futur &
      action située dans l'avenir\\[-.2mm]
    \textit{J'aurai fini à 18 h.} &
      \textit{Je finirai à 18 h.}
  \end{tabularx}
\end{whitebox}

\vspace{1.35mm}

% ------------------------------------------------------------
% SIGNAL WORDS
% ------------------------------------------------------------
\begin{goldbox}
  \begin{center}
    {\normalsize\bfseries\textcolor{Gold}{REPÈRES FRÉQUENTS DANS LE FUTUR}}
    \hspace{2mm}
    \signalpill{d'ici demain}\hspace{1mm}
    \signalpill{avant ce soir}\hspace{1mm}
    \signalpill{à midi}

    \vspace{0.75mm}

    \signalpill{d'ici vendredi}\hspace{1mm}
    \signalpill{quand}\hspace{1mm}
    \signalpill{lorsque}\hspace{1mm}
    \signalpill{dès que}
  \end{center}
\end{goldbox}

\end{document}
